\begin{enumerate}
\item Both $\angle ADB$ and $\angle AEB$ are right, so $D,E$ lie on the circle with diameter $AB$.
\item Both $\angle HDC$ and $\angle HEC$ are right, so $D,E$ lie on the circle with diameter $HC$.
\item In the latter circle $\angle HCD=\angle HED$. Since $E,H,B$ are on the same ray from $E$, $\angle HED=\angle BED$. In the first circle $\angle BED=\angle BAD=\angle HAB$.
\item Since $AH\perp BC$, $\angle HAB=90^\circ-\angle ABC$. Thus $CH$ makes the complementary angle with $CB$, giving $CH\perp AB$. This is exactly $CF\perp AB$.
\item The lines $AC,BC,HC$ are perpendicular respectively to $BH,AH,AB$. They are therefore the three altitudes of triangle $ABH$, meeting at $\boxed C$.
\item In triangle $AHB$, its two angles at $A,B$ are $90^\circ-B$ and $90^\circ-A$. Thus $\angle AHB=A+B=180^\circ-C$, proving the sum identity. Nevertheless $A,H,B,C$ are not concyclic: $H$ is strictly inside the acute triangle $ABC$, whereas a point on its circumcircle cannot be inside that triangle. The supplementary-angle criterion concerns opposite vertices of a simple convex cyclic quadrilateral; it cannot be applied in this concave boundary order.
\end{enumerate}
\begin{asy}
texpreamble("\usepackage{amsmath,amssymb}");
size(240);
defaultpen(fontsize(9pt));
pair A=(1,4),B=(0,0),C=(6,0),D=(1,0),E=(96.0/41,120.0/41),H=(1,1.25),F=(6.0/17,24.0/17);draw(circle((A+B)/2,abs(A-B)/2),lightgray);draw(circle((H+C)/2,abs(H-C)/2),lightgray);draw(A--B--C--cycle);draw(A--D,blue);draw(B--E,blue);draw(C--F,red);dot(H);label("$A$",A,N);label("$B$",B,SW);label("$C$",C,SE);label("$D$",D,S);label("$E$",E,(1,0));label("$F$",F,W);label("$H$",H,SE);
\end{asy}
